%!TEX root = ../../csuthesis_main.tex
\chapter{结论与展望}

本文以基于闭式解算法的持续学习方法及其实现为主题，探讨闭式解方法对于不同持续学习问题进行统一建模、推导及实验验证工作。为了解决传统的持续学习方法大多采用梯度下降方式导致容易发生灾难性遗忘并且回放机制存在隐私考虑和消耗大量空间资源的问题，本文从解析持续学习的角度出发，提出了用闭式解来代替增量过程中的梯度更新的方法，并将此应用于分类、分割、时间序列以及多模态大语言模型等多种任务中。在第六章第一节总结全文主要工作并回顾第三章统一数学框架与第四章四种类型的任务应用方法之间的联系，在第二节讨论该方法所存在的不足之处，在第三节提出未来的研究方向，在第四节进行总结。

\section{全文总结}

\subsection{统一数学框架的建立}

第三章给出了递归岭回归闭式解所必需的一般框架。即把连续学习中的固定编码器和更新分类头两部分整合成一个带有$\ell_2$正则化最小二乘问题，同时用逆自相关矩阵$\Psi_t$以及互相关矩阵$\mathbf{Q}_t$表示历史数据的信息，最后利用Woodbury矩阵恒等式得到仅需上一步骤的相关统计量以及当前步骤的新样本即可实现精确更新的方法。并且证明了递归更新的结果与对应的联合训练解析解是相同的，也就说回答了增量过程能否无需使用梯度进行更新的问题。

其一，从理论上讲，这一体系将稳定性和可塑性的冲突变为对二阶统计量进行代数上的管理和维护的问题，在相同理论上可以应用于不同的任务；其二，从算法角度看，每一阶段的学习难度从多轮反向传播降低到一次前向传播，避免了由于梯度下降导致的小范围最优值、超参数调整困难以及难以收敛等缺点；其三，只用存储大小为$d\times d$的数据统计量矩阵而不与历史样本数量相关联因此更适合长时间在线学习的情况。

\subsection{四类典型任务的应用扩展}

第四章基于第三章统一数学基础之上，提出了针对四种不同任务类型的闭式解持续学习的应用。对于医学图像疾病分类任务，将闭式解方法与冻结编码器、解析对齐、递归更新这三个步骤结合在一起，适用于隐私性强并且类别不断增加的医疗环境。对于时间序列分类任务，提出一种时序闭式解的方法，在此基础上采用多尺度特征融合、随机高维映射和集成扩展相结合的方式来解决灾难性遗忘和类内变化的问题。对于语义分割任务，提出一个针对二维图像以及三维点云的通用闭式解持续分割方法，利用冻结编码器、随机高维映射以及为这两者单独设计的伪标签，缓解语义漂移的问题。而对于多模态大语言模型专家路由任务，将闭式解的思想应用到了专家路由中，设计了一个解析式的任务选择头，实现了无需梯度就可以训练出路由器并且不需要存储任何历史多模态样本即可完成高效的专家路由的目的。


\subsection{统一实验体系下的系统验证}

第五章在统一实验框架下对以上方法进行了大量验证，包括医学图像分类、时间序列分类、二维图像及三维点云语义分割以及多模态大语言模型专家路由等四种类型的任务、共计十几个公开的数据集，均获得了比传统的持续学习基线更高的平均准确率以及更低的遗忘率。

相关研究成果已分别发表于 ACM International Conference on Multimedia（ACM MM，CCF-A）、IEEE International Conference on Acoustics, Speech and Signal Processing（ICASSP，CCF-B）以及 International Conference on Learning Representations Workshop（ICLR Workshop），另有 1 篇与时间序列持续学习相关的工作投稿至 Knowledge-Based Systems（KBS），当前在审稿阶段；对应论文均已在第四章各节脚注中给出标注。

\subsection{方法整体优势}

把以上几点结合起来看，基于闭式解的持续学习相比目前主流基于梯度下降的持续学习算法有四个方面优点并且与第一章四个主要成果对应。第一，在更新方式上，对应贡献一（提出使用闭式解代替增量阶段梯度更新的思想）：参数更新是通过一个明确表达式的直接给出，新任务对于旧任务的影响仅体现在对$\Psi_t$和$\widehat{\Phi}_t$的简单相加之上，从而避免了增量阶段不同任务梯度之间的冲突，从根本上杜绝了灾难性遗忘的可能性，而且递归解析解和联合训练解析解是等价的。第二，在整体框架上，对应贡献二（构建递归岭回归闭式解的整体数学框架）：该方法用一个统一的统计量维护的方式，使得各种各样的持续学习问题都可以在一个统一的理论下进行描述和解决。第三，在任务推广上，对应贡献三（将整个框架推广到四种常见任务并进行评估）：该方法具有较好的普适性，可以很自然地从标准类增量分类推广到稠密预测二维图像分割和三维点云分割、时间序列分类以及多模态大语言模型专家路由等，涵盖了这四类差别很大的任务，形成了一种理论和多个应用的研究模式。第四，在实用价值上，对应贡献四（在效率、隐私以及工程可行性上的好处）：一般情况下增量阶段只需要一轮解析更新而不需要多轮梯度迭代。同时由于这种方法不需要任何历史样本回放，只需要少量的二阶统计量就可以保存所有的历史信息，因此不受样本数量影响，也更加符合医疗影像、边缘计算设备等隐私要求高且硬件资源有限的应用场景。

\section{方法局限性}

尽管本文的方法对多个任务都得到很好的实验结果，但是也不能忽视它的不足之处，在表示的学习、解析分类头模型的能力、任务假设存在一些问题。对于这些问题进行分析有利于了解该方法的应用范围并为后续的研究提供思路。

\subsection{对初始阶段编码器的依赖}

本文采用冻结编码器和递归更新解析分类头两步走方案。这个方案能不能成功实现，取决于最初编码器学到表示是否有足够泛化能力。如果最初任务涵盖种类较少、分布偏差大或者是样本数量不够，那么之后仅通过解析分类头是无法补充底层表示缺失部分。也就是说，该方法有效范围在初始表征已经足够丰富情况下，而对于表征匮乏情况，则仍需使用更好预训练或自适应方法来弥补。

\subsection{解析分类头的表达力上限}

在编码器之后，本文采用线性解析分类头并用随机高维映射解决线性不可分问题。优点是闭式求解可行且具有递归迭代性质；缺点是在一定程度上，对于类内具有丰富模态、类间有显著非线性的数据集，该模型存在一定的上限。集成和多尺度特征融合可以提高此上线但是还是有限的。

\subsection{任务结构假设较强}

对于多模态大语言模型专家路由这部分内容，可做出三个假设：任务有明确分界点、专家数量可数、每个任务样本都能确定唯一的专家编号，在终身学习中这三个条件未必都满足——任务通常是按一定顺序不断到来，它们之间有很多重叠之处，而且新的任务也不一定就要有一个新专家来处理它。因此目前的方法过于关注任务身份，如果要使这个系统更加接近现实世界的应用，则还需要进一步完善其路由机制。



\section{未来研究展望}

为了解决上述问题，在此可以从表征基础、解析分类头模型、任务结构等方面对今后的研究工作进行一些设想。

\subsection{结合大规模预训练与基础模型}

未来可以考虑用更大规模预训练模型或者基础模型作为初始编码器，如 DINOv2、SAM、CLIP 和多模态大语言模型等，从而提升底层表示迁移能力。一方面，预训练模型跨域任务上良好零样本性能大大减少对于初始任务覆盖类别的要求；另一方面，采用分阶段微调或持续预训练方法，使编码器在不影响先前学习结果情况下进行一定程度调整，是将静态预训练和解析持续学习两者结合起来有效手段。

\subsection{解析分类头模型的非线性扩展}

在线性可分时，本文已利用随机高维映射对其进行一定的增强；但是当环境变为更为复杂的真实环境后，模型所能学到的最大能力就达到了瓶颈。未来的研究可以朝向更加丰富的非线性解析建模。一种思路是引入随机傅里叶特征对模型进行核化扩展。

\subsection{面向开放世界的任务结构扩展}

为进一步改进目前基于多模态专家路由方法对于任务边界明确、专家可列举的前提条件所存在的问题，未来的研究可以考虑如下几个方面的工作：第一，将硬路由转化为软路由，使得该方法能够适用于任务之间存在语义重叠的情况；第二，在系统中加入OOD（开放类别）检测以及新的专家自动生成的能力，让系统可以在未知领域中也能具有可解释性和可追踪性的表现。


\section{本章小结}

本章对全文工作进行了总结并对其未来的研究方向进行了展望。本文是以第三章统一方法为基础，在此基础上对四种不同的任务进行研究，在第四章中提出相应的算法并进行实验验证，在第五章中利用统一的实验框架进行系统的比较。基本实现了基于闭式解持续学习的研究。本文的研究成果表明闭式解持续学习不仅是一个有效的减少遗忘的新途径同时也是一种有较高的理论性和实践可行性的持续学习的方法。但是本论文所提出的算法也有其不足之处：如对于编码器的依赖性较强、解析模型的能力有限、任务设置过于简单以及数值不稳定等问题。因此后续的研究可以从以下几个方面入手：一是结合大规模预训练基础的模型；二是对解析分类头模型进行非线性扩展；三是开放世界的任务设置等。希望本文的研究工作可以给持续学习领域提供一种新的思路并且也为闭式解的发展指明方向。